% ============================================================
% paper/appendices/appx_bibliographic_positioning.tex
% Appendix B: Bibliographic Positioning (Non-Comparative)
% Purpose: situate conceptual lineages without validation claims.
% ============================================================

\section*{Appendix B: Bibliographic Positioning (Non-Comparative)}

\subsection*{Purpose of This Appendix}

This appendix reduces bibliographic isolation by placing the
Enterprise Continuum $K_{\mathrm{EA}}$ in a set of
\emph{conceptual lineages} that use structurally related notions
(state spaces, constraints, thresholds, invariants, termination).
It does \textbf{not} validate, compare, or evaluate this whitepaper
against any external literature.

The appendix is explicitly \emph{non-exhaustive}. Its function is
reader orientation only.

\subsection*{Interpretation Boundary (Non-Claims)}

The presence of related conceptual lineages does not imply that:
(i) this whitepaper is ``based on'' any specific prior framework,
(ii) any cited work endorses the present formal object, or
(iii) any empirical, predictive, or practical claim follows.

This document remains an executable-closed projection of
\texttt{ETS.K\_EA.v1.1}. No external reference has normative authority
over the ETS.

\subsection*{Related Conceptual Lineages (Orientation Only)}

The following lineages are structurally adjacent in the sense that they
provide vocabulary for \emph{closed formal objects} and their
\emph{admissibility/termination boundaries}:

\begin{itemize}
  \item \textbf{Formal systems, specification, and invariants.}
  Axiomatic description, invariants, and No-Go style constraints appear
  broadly in formal methods and specification disciplines.

  \item \textbf{Dynamical systems and state-space views.}
  The practice of describing systems by a set of admissible states and
  constraints is common in dynamical systems, control, and systems
  engineering.

  \item \textbf{Thresholds, phase transitions, and admissibility.}
  The use of explicit thresholds separating admissible from
  non-admissible regimes has analogues in physical phase transition
  formalisms and in constraint-driven regime changes.

  \item \textbf{Cybernetics, feedback, and closure motifs.}
  The emphasis on cycles/closure as structural conditions resonates with
  classical feedback-oriented treatments in cybernetics and systems
  theory.

  \item \textbf{Enterprise architecture as structural description.}
  EA traditions provide long-standing efforts to describe enterprises as
  structured systems (models, viewpoints, constraints), without requiring
  any shared claim about effectiveness.

  \item \textbf{Reproducibility and artifact identity.}
  The use of hashes and immutable references aligns with modern
  reproducibility norms in software and computational research, where
  identity is anchored by frozen objects rather than interpretive text.
\end{itemize}

\subsection*{Representative References (Non-Exhaustive)}

The list below is intentionally minimal and representative; it is not a
``related work'' section.

\begin{itemize}
  \item \textbf{Formal methods / specification:}
  Hoare (1969) on axiomatic reasoning about programs;
  Lamport (1994) on formal specification; ISO/IEC/IEEE 42010 (systems and
  software architecture description.\footnote{C.\ A.\ R.\ Hoare,
  ``An Axiomatic Basis for Computer Programming,'' \emph{Communications of
  the ACM}, 1969.}
  \footnote{L.\ Lamport, \emph{Specifying Systems}, Addison-Wesley, 1994.}
  \footnote{ISO/IEC/IEEE 42010:2011, ``Systems and software engineering ---
  Architecture description.''}

  \item \textbf{Dynamical systems / control:}
  standard state-space formulations in systems and control
  (e.g., Kalman-style representations).\footnote{R.\ E.\ Kalman,
  ``On the General Theory of Control Systems,'' 1960.}

  \item \textbf{Thresholds / phase transitions:}
  canonical expositions of critical phenomena and regime separation.\footnote{
  N.\ Goldenfeld, \emph{Lectures on Phase Transitions and the Renormalization
  Group}, Addison-Wesley, 1992.}

  \item \textbf{Cybernetics / feedback:}
  closure and feedback as structural motifs.\footnote{
  N.\ Wiener, \emph{Cybernetics: Or Control and Communication in the Animal
  and the Machine}, MIT Press (reprint), 1961.}

  \item \textbf{Enterprise architecture (descriptive traditions):}
  long-standing architecture description and modeling conventions.\footnote{
  The Open Group, \emph{TOGAF Standard} (architecture framework), various
  editions.}

  \item \textbf{Artifact identity / reproducibility:}
  content-addressing and integrity anchoring as identity mechanisms.\footnote{
  S.\ Chacon and B.\ Straub, \emph{Pro Git}, 2nd ed., Apress, 2014
  (hash-based object identity in Git).}
\end{itemize}

\subsection*{Final Note}

This appendix provides \emph{orientation} only.
No inference about empirical validity, effectiveness, or utility is
licensed by the existence of these lineages.

%%%End of Document%%%
